\section{Organisation and implementation of the project}

%This paragraph refers to the evaluation criterion 
%``Organisation and implementation of the project'' whose 
%sub-criteria are differentiated according to the funding instruments. 

%The choice of the funding instrument requires reading the expectations 
%in terms of objectives and partnership or team of each of the instruments 
%before making this choice. This choice is final throughout the AAPG selection process.
 
\projacro\ is a PRME project that will be hosted at Aix-Marseille Université (amU), 
particularly at the Centre de Physique des Particules de Marseille (CPPM). 
The team consists of researchers already active in the field, all members of either 
DESI and ZTF collaborations, spaning a wide range of expertise. 

%\subsection{The scientific coordinator: Julian Bautista (90\% ETPR)} 



%describe the scientific coordinator, her/his experience in the scientific
    %field, her/his involvement in the project including her/his involvement rate;
    
\subsection{The scientific team at the CPPM}


     % Describe the team, its expertise and complementarity to achieve the 
     %scientific goals, identity of the scientists involved and their involvement rate,
     % demonstration of the sustainability of the team throughout the project.

     \begin{enumerate}
     \item \textbf{Julian Bautista - Scientific Coordinator (Professor, 100\% ETPR)} 
     
     Expert in galaxy clustering, 
    type-Ia supernovae and peculiar velocities for cosmology. 
    Major contributor of past cosmological results with spectroscopic surveys 
    such as SDSS-III BOSS, SDSS-IV eBOSS.
    Currenly active member of DESI and ZTF and actively working 
    in peculiar velocities for the past 5 years, with the support of the 
    Chaire d'Excellence A*Midex.
    Current co-chair of the DESI Transients and Low-redshift Cosmology Working Group, 
    whose main goal is to perform growth rate measurements with the DESI Peculiar Velocity Survey. 
    
        
    %\item \textbf{Fei Qin:} 
    %    Postdoctoral researcher in the group of J. Bautista. He has more than four years
    %    of experience in cosmology with peculiar velocities. He first-authored more than 
    %%    six publications in the field, including measurements of the local cosmic bulk flow, 
    %    of the growth rate of structures with several past peculiar velocity surveys with 
    %    the momentum power spectrum technique. He also has valuable expertise in the 
    %    production of synthetic datasets. He is currently one of the main authors of the DESI PV 
    %    DR1 results. His expertise will be key in the project, particularly with the supervision of
    %    future PhD students. His current contract ends mid 2026 and we wish to extend his stay
    %    with the funds of the project. 

    %\item \textbf{Damiano Rosselli:} 
    %    Recently graduated PhD student, he is currently a postdoctoral researcher in the group of J. Bautista. 
    %    D. Rosselli is an expert in the measurement of the growth rate of structures 
    %    with the maximum likelihood estimator method, 
    %    and he focused his past research on realistic measurements of the growth rate of structures 
    %    with synthetic samples of LSST supernovae, which he created himself \cite{rosselliForecastGrowthrateMeasurement2025}.
    %    D. Rosselli has expertise in the study of observational systematics and how they impact 
    %    cosmological results, and his knowledge will be key for the success of the \projacro. 
    %    His contract also ends mid 2026. 

\item   \textbf{Dominique Fouchez (CNRS Research Director, 10\% ETPR)} 
        
Expert in supernova cosmology (SNLS, ZTF and LSST), 
photometric detection of sources as well 
as applications of machine learning to astronomy. 
Currently PhD supervisor of Rafael Kebadian (see below).
    
\item \textbf{Benjamin Racine (CNRS Researcher, 10\% ETPR)}

Expert in cosmology with the CMB and SNIa.  
Leading the efforts to homogenize the flux calibration of ZTF images 
to produce datasets suited for precision cosmology with SNIa. 
His raw data expertise is of extreme value for \projacro, as he will be able to advice on 
potential sources of systematic effects that impact cosmological results.  
Currently PhD co-supervisor of Rafael Kebadian (see below).

\item  \textbf{Fabrice Feinstein (Professor, 40\% ETPR)} 
    
Expert on ZTF datasets. He has been closely working with B. Racine on
on the flux calibration of
ZTF image and has experience in supernova cosmology. Will contribute 
with the investigation of observational systematics in the data (WP~\ref{systematics}). 

\item  \textbf{Juan Mena-Fernandez (Postdoc, 30\% ETPR)}

Expert in galaxy clustering, weak lensing and mock catalog production. 
He is active member of DESI, DES and Rubin-LSST. 
His expertise will be useful for WP~\ref{sims} and WP~\ref{cosmo}. 

\item \textbf{Corentin Hanser (Postdoc, 30\% ETPR)}
    
Transitioning from cosmology with galaxy clusters to peculiar velocities.  
He has been active in setting up a careful blinding procedure for 
growth rate measurements with SNIa peculiar velocities. 
He will contribute during one year to the project, in particular to WP1. 

\item \textbf{Rafael Kebadian (PhD student, 30\% ETPR)}
     
Currently leading the measurement of the growth rate of structures 
with half of ZTF data (DR2). He will be part of the
project during his last year and will contribute to the first 
two WP of \projacro.   
        

     \end{enumerate}


\small
\begin{center}
  \textbf{Implication of the scientific coordinator in on-going projects \\ \ \\}

\begin{tblr}{
    width=\linewidth,
    colspec={X[l,valign=m,wd=2.5cm]
             X[l,valign=m,wd=1.5cm]
             X[l,valign=m,wd=2.5cm]
             X[l,valign=m,wd=2.cm]
             X[l,valign=m]
             X[l,valign=m,wd=2cm]},
    row{1}={font=\bfseries},
     hline{2} = {solid,gray}
}
\hline
\hline 

Name of Researcher &
Person. month &
Call, funding agency, grant allocated &
Title of the project & 
Name of the scientific coordinator& 
Start - End \\


Julian Bautista & N/A & N/A & N/A & N/A & N/A \\ 

\hline 
\hline

\end{tblr}
\end{center}


\subsection{Implemented and requested resources to reach the objectives}

Including fees, \textbf{the total budget of \projacro\ amounts to 
473 697,93~€ for a duration of 48 months.} 
It is divided between staff expenses, travel support, computers, conference fees 
and conference organisation, plus admin fees. 


%\partnerTitle{Partner 1: XXX}


\ressourceTitle{Staff expenses: 391,3k~€}

This amount would allow us to hire:
\begin{itemize}
    \item {one PhD student for 36 months for a cost of 137 977,00~€};
    \item {one senior postdoctoral researcher for 36 months for a cost of 253 378,00~€}. 
\end{itemize}

The senior postdoctoral researcher would have more than 2 years of experience after their PhD degree. 

\ressourceTitle{Equipments: 6k €}

Two personnal computers for the future PhD student and postdoctoral 
researcher, each one at the maximum cost of 3k~€. Recently the price 
of RAM memory increased considerably, driving up the cost of new computers. 

Those computers will be essential to the project as they will serve 
the two new members of the team to develop analysis code and acces 
the supercomputers (Mardec, NERSC) where the high-performance computing 
will be performed. 




\ressourceTitle{Service delivery: 3 k€}

We will publish our findings in the Astronomy \& Astrophysics  journal (A\&A).
Recent changes to publication policy limit the number of pages to a maximum of 12
without publication fees. Currently, each extra page costs 150~€. Since our 
publications often exceed this limit, we will request 3k~€ for publication fees. 

As a reminder, open access versions of our published articles will be posted on
the arXiv website without any associated fees. 


\ressourceTitle{Other costs: 30,5 k€}


We would like to ask \textbf{19k~€ for travel support} for each member 
of the team, including the scientific coordinator. 
Each new member of the team will be assisted with 2,5k~€ per year ($2\times 3 \times 2,5 = 15k~€$), 
while the scientific coordinator will be assisted with 1k~€ per year ($4 \times 1k~€ = 4k~€$). 

We ask an additional \textbf{4k~€ for the payment of conference fees} for the same people
for the duration of the project. 

To increase the visibility of \projacro, as well as our local institutions, 
we would like request \textbf{7,5k~€ to organise a version of the CosmicFlows conference in Marseille}. 
Such funds will help with the costs of the conference, such as room rentals, 
coffee breaks and conference dinner. We expect to welcome around 50 people 
for a duration of a week, to share about the domain of cosmology and astrophysics 
with peculiar velocities. 



%\smallSection{Requested means by item of expenditure and by partner}


% \begin{center}


\begin{center}

    \textbf{Requested means by item of expenditure and by partner }

    \begin{tblr}{
    colspec={lr},
    row{1}={font=\bfseries},
    }
    \hline
    \hline 

    \textbf{Type of expenditure}
    &
    {\centering Amount (€) } \\
    
    \hline

    {Staff expenses} &
    391 355,00  \\
    
    {Equipments } &
    6~000,00 \\
    
    {Service delivery} &
    3~000,00  \\
    
    {Other costs} &
    30~500,00  \\
    
    {Structure costs**} &
    58~165,43  \\
    \hline    

    \textbf{Sub-total} &
    489 020,43  \\
    \hline
    \textbf{Requested funding} &
    489 020,43  \\
    \hline 
    \hline 
    
\end{tblr}
\end{center}


